% Copyright (c) 2026 by C. Klukas.
% Licensed under the MIT License. See LICENSE for details.

A paglet reaches another paglet only through an endpoint, a capability
that names the target by its ID (Section~\ref{sec:model}). The endpoint
does not name a host. This chapter shows how a message still arrives
when its target has moved since the endpoint was handed out, what the
sender sees when it cannot, and what holds a paglet in place while
somebody acts on an answer.

\subsection{Endpoints across hosts}\label{sec:communication-endpoints}

The host of a sending paglet keeps, with every endpoint and reply
capability whose target is on another host, a \emph{host hint}: the key
ID of the host where the target was last known to be (\code{Cap::host} in
the runtime's \code{capability.hpp}; empty for targets on this host).
A capability gets the hint when it leaves a host, in a message or with a
moving paglet. Endpoints to system paglets carry none: they do not travel
in messages, and a moving paglet's endpoints to its default services are
made anew on arrival from the destination's own system paglets
(Section~\ref{sec:movement}). A send, request or reply whose target is not
on this host leaves the runtime as a \code{RemoteMessage} (the runtime's
\code{mobility.hpp}): the target and its hinted host, the message, and
for a request the requester, its host and the correlation. Only endpoints
to paglets travel between hosts, and artifact capabilities (of the
\code{artifacts} system paglet, Section~\ref{sec:mesh}), which name the
host that holds their content. A send, request or reply to a paglet on
another host that carries any other capability is refused with
\code{unsupported}; a message that was accepted on the same host and then
follows its target after a move has such capabilities dropped as
\code{gone}. Lending a capability (\code{SendOptions::lend}) shows it to a
system paglet of this host only; a lend to another host is
\code{invalid\_argument}. The node wraps the message in a \code{deliver}
frame and sends it through the channel to the hinted host
(Section~\ref{sec:communication-channel}).

The hint is a guess. It is right for a target that has not moved. When
it is wrong, the host that receives the frame has to find the target
without help from the sender.

\subsection{Location records}\label{sec:communication-records}

Every paglet has a location record: the host that holds it, its move
counter, when it arrived there and its owner. No host is special. The
records of a paglet live on three \emph{responsible hosts} chosen by
consistent hashing~\cite{karger1997consistent} over the hosts this host
believes are up: each host
has 16 points on a 64-bit ring (SHA-256 over a fixed prefix, the host key
and the point number), a paglet has one point, and the first three
distinct hosts clockwise from it keep its record. With three or fewer
hosts, all of them do (Listing~\ref{lst:comm-responsible}).

% snippet: cpp/host/src/node/location.cpp#responsible [float,caption={The responsible hosts of a paglet: the first three distinct hosts clockwise from its point on the ring},label={lst:comm-responsible}]
\begin{lstlisting}[style=paglets-cpp,float,caption={The responsible hosts of a paglet: the first three distinct hosts clockwise from its point on the ring \src{cpp/host/src/node/location.cpp}},label={lst:comm-responsible}]
std::vector<mesh::PublicKey> Node::Impl::responsible_for(const runtime::PagletId& paglet) {
    update_view();
    if (view.size() <= replicas) return view;
    std::vector<mesh::PublicKey> out;
    const std::uint64_t p = paglet_point(paglet);
    auto it = std::ranges::lower_bound(ring, std::pair{p, mesh::PublicKey{}});
    for (std::size_t n = 0; n < ring.size() && out.size() < replicas; ++n, ++it) {
        if (it == ring.end()) it = ring.begin();
        if (std::ranges::find(out, it->second) == out.end()) out.push_back(it->second);
    }
    return out;
}
\end{lstlisting}
% /snippet

Creation writes a record with counter~0. Every move increments the
counter, and a record with a higher counter replaces one with a lower
counter; the same counter from the same host only refreshes it. The
counter is what makes a stale record harmless: a host that was down and
returns with an old record loses to any newer one. The update is the
commit step of the move protocol (Section~\ref{sec:movement-protocol}):
the source sends the new record to the responsible hosts and commits the
move only after a majority ($\lfloor n/2 \rfloor + 1$) acknowledged it,
recording above a newer record at most three times before it aborts and
the paglet stays. So every committed move is visible to a majority of
the responsible hosts before the paglet runs on its new host.

The set of live hosts changes. Every host sends a heartbeat to every
enrolled host every 2\,s (any frame counts as one); a host not heard from
for 10\,s, or one that announced another mesh protocol or ABI version, is
out of this host's view, and the ring is rebuilt over the view. Holders
then re-send the records of their paglets, and every host sends the
records it keeps to hosts that became responsible for them. Holders also
refresh their records every five minutes, and a record not refreshed for
an hour ends, which removes the records of paglets that ended. Records are
kept in the node's state directory, so a restarted host still has them.
Different hosts can disagree about the view for a few seconds; the lookup
below has a fallback for the case that no responsible host knows a paglet.

\subsection{Delivery after a move}\label{sec:communication-delivery}

A lookup (the node's lookup behind \code{Node::locate}, the
\code{locator} system paglet, \code{paglets-host remote locate} and
message forwarding) runs in three steps, the second and third limited to
3\,s each. First, a paglet on this host is found at once. Second, the
responsible hosts are asked (\code{loc-get}); with answers from a
majority, the record with the highest counter wins, which is the latest
committed location. When the 3\,s pass with fewer answers, the best record
that did arrive is used, which may be stale. Third, if no responsible host
knows the paglet, every live host is asked (\code{loc-who}), and the host
that holds it answers; if none does, the lookup ends with
\code{not\_found}. Hosts cache the answers of lookups and moves for 10\,s.

Message forwarding adds a cheaper first step. A host remembers for ten
minutes where a departed paglet went (a \emph{tombstone},
\code{Config::tombstone\_ttl}). A message that arrives for a paglet that
left is forwarded to the tombstone's host, with a hop counter in the
message that allows at most four forwards in all. Only when the host has
neither the paglet nor a tombstone, and the message has not yet been
forwarded four times, does the runtime hand the message to the node
(\code{MobilityHooks::unresolved}), which uses its location cache or
starts a lookup and forwards the message to the host it finds
(Listing~\ref{lst:comm-unresolved}). A request that nobody holds is
answered \code{not\_found} to the requester; a one-way message that
nobody holds is dropped. Figure~\ref{fig:comm-forward} shows both paths.

% snippet: cpp/host/src/node/location.cpp#unresolved [float,caption={Messages for paglets that are neither here nor remembered: the location cache, otherwise a lookup; requests for paglets nobody holds are answered},label={lst:comm-unresolved}]
\begin{lstlisting}[style=paglets-cpp,float,caption={Messages for paglets that are neither here nor remembered: the location cache, otherwise a lookup; requests for paglets nobody holds are answered \src{cpp/host/src/node/location.cpp}},label={lst:comm-unresolved}]
for (auto& m : queue) {
    auto deliver = [this](runtime::RemoteMessage& msg, const mesh::PublicKey& host) {
        msg.host = mesh::key_id(host);
        ++msg.hops;
        send_remote(std::move(msg));
    };
    if (auto c = loc_cache.find(m.target); c != loc_cache.end() && Clock::now() - c->second.second < cache_ttl &&
                                           c->second.first.host != key.public_key()) {
        deliver(m, c->second.first.host);
        continue;
    }
    Lookup l;
    l.kind = LookupKind::locate;
    l.paglet = m.target;
    l.done = [this, m = std::move(m), deliver](const LookupResult& r) mutable {
        if (r.status == abi::ok) return deliver(m, r.where.host);
        if (m.kind != runtime::RemoteMessage::Kind::request) return;
        // Nobody has it: the requester hears so.
        runtime::RemoteMessage answer;
        answer.kind = runtime::RemoteMessage::Kind::reply;
        answer.target = m.requester;
        answer.host = m.requester_host;
        answer.correlation = m.correlation;
        answer.status = abi::not_found;
        send_remote(std::move(answer));
    };
    start_lookup(std::move(l));
}
\end{lstlisting}
% /snippet

\begin{figure}[htbp]
\centering
\begin{tikzpicture}[x=1cm,y=1cm]
  \node[box] (s) at (0,0) {sender's host $S$};
  \node[box] (a) at (3.9,0) {old host $A$};
  \node[box] (r) at (7.8,0) {responsible host $R$};
  \node[box] (b) at (11.7,0) {new host $B$};
  \foreach \n in {s,a,r,b} \draw[lifeline] (\n.south) -- ++(0,-5.0);
  \draw[msg] ($(s.south)+(0,-0.6)$) -- node[above]{\code{deliver} (hint: $A$)} ($(a.south)+(0,-0.6)$);
  \draw[msg] ($(a.south)+(0,-1.3)$) -- node[above,pos=0.3]{tombstone for $B$: \code{deliver}, hop $1$} ($(b.south)+(0,-1.3)$);
  \draw[dashedarrow] ($(a.south)+(0,-2.3)$) -- node[above,event]{no tombstone: \code{loc-get}} ($(r.south)+(0,-2.3)$);
  \draw[dashedarrow] ($(r.south)+(0,-2.9)$) -- node[above,event]{\code{loc-rec}: $B$, counter $c$} ($(a.south)+(0,-2.9)$);
  \draw[dashedarrow] ($(a.south)+(0,-3.6)$) -- node[above,event,pos=0.3]{\code{deliver}, hop $1$} ($(b.south)+(0,-3.6)$);
  \draw[msg] ($(b.south)+(0,-4.5)$) -- node[above,pos=0.7]{reply (requester's host: $S$)} ($(s.south)+(0,-4.5)$);
\end{tikzpicture}
\caption{A request to a paglet that moved from $A$ to $B$ after the sender got its endpoint. $A$ forwards by its tombstone (solid), or, after the tombstone ended or $A$ restarted, by a lookup on the responsible hosts (dashed). The reply goes to the requester's host directly, because a request names it.}
\label{fig:comm-forward}
\end{figure}

A paglet in the middle of a move handles no message: after the handler
that called \code{dispatch} it receives only the \code{dispatching} event,
and from then until the outcome is known the runtime holds the messages
that arrive for it. When the move succeeds, the held messages are sent on
to the new host in mailbox order (by priority, then arrival); a held
request becomes a remote request that names the requester and its host,
so its reply finds the way back. When the move fails, the paglet
continues on its host and handles the held messages there. In the other
direction, the paglet's own requests in flight travel with it: the
travelling state lists each pending correlation with its remaining time,
and the new host arms the timeouts again on arrival. The continuation of
a request lives in the paglet's memory and moves with the image, so a
reply that arrives after the move runs the continuation on the new host.
Table~\ref{tab:comm-timing} collects the constants of this section.

\begin{table}[htbp]
\centering\small
\caption{Timing and limits of location and delivery. Names without a file are fields of \code{Node::LocationTiming} (\code{node.hpp}). All are compiled-in defaults; the node's timings can be changed by its embedder, not from the command line.}
\label{tab:comm-timing}
\begin{tabular}{llll}
\toprule
Quantity & Value & Name & File \\
\midrule
Heartbeat to every enrolled host & 2\,s & \code{heartbeat} & \\
A host counts as down after & 10\,s & \code{host\_timeout} & \\
Lookup steps two and three & 3\,s each & \code{lookup\_timeout} & \\
Location cache & 10\,s & \code{cache\_ttl} & \code{location.cpp} \\
Record refresh by holders & 5\,min & \code{refresh} & \\
A record not refreshed ends after & 1\,h & \code{record\_ttl} & \\
Tombstone of a departed paglet & 10\,min & \code{Config::tombstone\_ttl} & \code{runtime.hpp} \\
Forwards of one message & 4 & \code{max\_forwards} & \code{runtime.cpp} \\
Redirects a pin follows & 8 & \code{max\_hops} & \code{location.cpp} \\
Waits for a moving paglet & $40 \times 100$\,ms & \code{max\_busy}, \code{busy\_wait} & \code{location.cpp} \\
Default request timeout & 30\,s & \code{abi::default\_timeout\_ms} & \code{abi.hpp} \\
Mailbox of a paglet & 1024 messages & \code{Config::mailbox\_limit} & \code{runtime.hpp} \\
\bottomrule
\end{tabular}
\end{table}

\subsection{Pins}\label{sec:communication-pins}

An answer from the locator can be stale by the time its receiver acts on
it. A \emph{pin} keeps a paglet on its host for a while. While a paglet
holds an unexpired pin, its own \code{dispatch} returns \code{pinned}
(error $-17$), and a dispatch an owner or admin requests is refused with
the same error. A dispatch requested before the pin fails with
\code{paglets.move\_failed} when it comes due.
\code{self\_info().pinned\_until} names the end of its last pin. Cloning,
deactivation and messages are not affected. A paglet can hold up to 64
pins and moves once all have ended. Pins are leases: they end at their
expiry, when their holder releases them or when an admin ends them
(\code{remote unpin}). The node stores them in its state and applies them
again after a restart. A pin never races a move: pinning a paglet that is
leaving or about to leave is refused with \code{bad\_state}, and the pin
protocol then waits 100\,ms and looks the paglet up again; a host that no
longer has the paglet redirects the pin to its tombstone's host. One
counter bounds both: a pin ends with \code{timeout} after eight redirects,
or after 40 waits and redirects together.

Paglets pin through the \code{locator} system paglet, whose endpoint
every paglet has by default (operations \code{locate},
\code{locate\_and\_pin} and \code{release}). The caller lends the
endpoint of the target, which must carry the \code{locate} or \code{pin}
right. \code{locate\_and\_pin} finds the paglet and pins it where it is in
one step; its reply carries the \code{Location} (paglet ID, host key ID
and name, time of the last move, move counter), the pin's ID and end, and
as its first capability a \emph{pin capability} of kind \code{pin} with
the right \code{release}. The capability is transferable and stays valid
when its holder moves; releasing revokes it and every copy. Who may pin,
and for how long, is decided by the mesh policy: \code{locate\_and\_pin}
needs a rule that allows service \code{locator}, operation \code{pin},
for the caller, and the rule's \code{max\_duration\_ms} caps the pin (one
hour without a maximum; the request itself is clamped to between 1\,s and
30 days). Without such a rule, pinning is \code{denied}. Owners (for their
own paglets) and admins pin through a CLI session (\code{paglets-host
remote pin}), which needs no rule: 10 minutes by default, at most 24
hours.

The Hide and Seek example (\code{cpp/examples/seek/}) plays this out
under continuous movement. The hider is a child that moves to \code{any}
host on a timer; while pinned, its \code{dispatch} says so at once and it
tries again later (Listing~\ref{lst:comm-hider}).

% snippet: cpp/examples/seek/seek.cpp#hider [float,caption={The hider moves on a timer and treats a refused dispatch as a reason to try later},label={lst:comm-hider}]
\begin{lstlisting}[style=paglets-cpp,float,caption={The hider moves on a timer and treats a refused dispatch as a reason to try later \src{cpp/examples/seek/seek.cpp}},label={lst:comm-hider}]
    void on_created(paglets::Started& s) override {
        // Created by a seeker: hide (args: the hop interval).
        if (!s.decode_args(hop_ms_) || hop_ms_ <= 0) return;
        hiding_ = true;
        (void)paglets::after_raw(hop_ms_, "hop", {});
    }

    void on_arrived(const paglets::abi::ArrivedEvent&) override {
        if (!hiding_) return;
        ++moves_;
        (void)paglets::after_raw(hop_ms_, "hop", {});
    }

    void on_move_failed(const paglets::abi::MoveFailed&) override {
        if (hiding_) (void)paglets::after_raw(hop_ms_, "hop", {});  // pinned, or nowhere to go: later
    }

private:
    // -- the hider ----------------------------------------------------------------------

    void hop() {
        if (!hiding_) return;
        // While pinned, dispatch says so at once: try again later.
        if (!paglets::dispatch("any")) (void)paglets::after_raw(hop_ms_, "hop", {});
    }
\end{lstlisting}
% /snippet

The seeker keeps one endpoint to the hider. Each round it pins the hider
for at most 10\,s with \code{patterns::with\_pinned}
(\code{<paglets/patterns/locate.hpp>}), which calls
\code{locate\_and\_pin}, runs the work with the location, and releases
the pin when the work calls \code{done}. Pinned, the hider answers the
request \code{moves} from the host where it was found, which is the point
of the pin: the answer and the location belong to the same host
(Listing~\ref{lst:comm-seek}). A round starts only once the hider has
moved since the last catch, so every catch finds it after a new move.
Table~\ref{tab:comm-seek} records the test on three hosts in one process,
with a hop every 30\,ms.

% snippet: cpp/examples/seek/seek.cpp#seek [float,caption={A round of Hide and Seek: pin the hider where it is, ask it something there, let it go},label={lst:comm-seek}]
\begin{lstlisting}[style=paglets-cpp,float,caption={A round of Hide and Seek: pin the hider where it is, ask it something there, let it go \src{cpp/examples/seek/seek.cpp}},label={lst:comm-seek}]
// The next round starts once the hider has moved since the last catch,
// however slow the hosts are.
void round() {
    if (static_cast<std::int64_t>(score_.catches.size()) >= game_.rounds) return finish();
    hider_moves([this](std::int64_t n) {
        if (!score_.catches.empty() && n >= 0 && n <= caught_at_moves_) {
            (void)paglets::after_raw(game_.pause_ms, "next-round", {});
            return;
        }
        seek();
    });
}

void seek() {
    pt::with_pinned(hider_, 10'000, "hide and seek",
                    [this](paglets::Result<pt::Location> where, std::function<void()> done) {
                        Catch c;
                        if (where) {
                            c.host_name = where->host_name;
                            c.moves = where->moves;
                        } else {
                            c.error = std::string(paglets::abi::error_name(where.error()));
                        }
                        score_.catches.push_back(std::move(c));
                        // Pinned, it stays where it is while it counts.
                        hider_moves([this, done](std::int64_t n) {
                            caught_at_moves_ = n;
                            done();  // let it go
                            (void)paglets::after_raw(game_.pause_ms, "next-round", {});
                        });
                    });
}
\end{lstlisting}
% /snippet

% Hand-written from docs/results/seek.json (recorded 2026-10-12): the
% `% results: seek` marker renders every column of the JSON (test name,
% outcome, platform, date, ...), which is wider than the text block; see
% the refinement plan. Keep these numbers in step with the JSON.
\begin{table}[htbp]
\centering\small
\caption{Hide and Seek on three hosts in one test process (\code{example: Hide and Seek finds and pins a paglet that keeps moving}, Darwin arm64, Release build, recorded 2026-10-12): four rounds with a hop every 30\,ms and 300\,ms between rounds. Both runs passed.}
\label{tab:comm-seek}
\begin{tabular}{lrrr}
\toprule
Mode & Duration (s) & Catches & Moves of the hider \\
\midrule
in-process & 1.7 & 4 & 35 \\
worker processes & 1.8 & 4 & 30 \\
\bottomrule
\end{tabular}
\end{table}

\subsection{Asynchronous sends, replies and failures}\label{sec:communication-async}

\code{Endpoint::send} and \code{Message::reply} are synchronous by
default: the call returns when the host has accepted the message, with
\code{denied}, \code{expired} or \code{quota} when the endpoint does not
allow it, \code{not\_found} when a target on this host ended, and
\code{quota} when the target's mailbox here is full. For a target on
another host the call succeeds once the runtime has handed the message to
the node; whether the frame reaches the other host is not reported. The
remote host answers a request that it cannot deliver (unknown paglet, full
mailbox, a system paglet of its own) with the error as the reply status,
and drops a one-way message it cannot deliver.

With \code{SendOptions::async}, \code{reply\_async} or \code{async = true}
in \code{reply\_to}, the call returns \code{ok} as soon as the document is
well-formed. This matters in worker processes (Section~\ref{sec:model}):
a synchronous call is a round trip to the host, an asynchronous one is a
one-way frame. A failure is reported later as the host's report
\code{paglets.undelivered}, which the SDK delivers to
\code{on\_undelivered} with the name, the handle and the error the
synchronous call would have returned; transferred capabilities are gone
and a failed reply has answered the requester \code{gone}
(Section~\ref{sec:lifecycle-reports}). A request is never
asynchronous: \code{RequestOptions} has no \code{async}, and the ABI call
refuses the flag with \code{invalid\_argument}, since a request has its
continuation. Its timeout (\code{RequestOptions::timeout\_ms}, 30\,s by
default) runs on the requester's host and is armed again after a move;
when it passes, the continuation runs with status \code{timeout}, and a
reply that arrives later is discarded.

Ping Pong (Listings~\ref{lst:pingpong-roles} and \ref{lst:pingpong-bounce})
has both sides: the child answers with \code{reply\_async}, the parent
answers a deferred reply from the continuation of the last ping or drops
it, which answers the caller \code{gone} at once. Nothing in it depends on
where the child is: had the child dispatched itself between two pings,
the parent's endpoint would carry the next ping to it by the paths of
Section~\ref{sec:communication-delivery}. The example in the repository
stays on one host; Hide and Seek above and the Latency Map below talk
across hosts.

\subsection{Hosts without an inbound port}\label{sec:communication-relays}

A host behind NAT or a firewall can open connections but accept none.
Started with \code{--no-listen}, such a host keeps two \emph{relays},
reachable hosts that it polls for the frames held for it (the
transport's \code{poll} request), and announces them to the mesh in place
of an address; the bootstrap contact of \code{--join} is the first. Every
other host reaches it through a relay over a Noise channel of its own,
tunneled in transport control frames, so the relay sees only the two host
keys and the sizes. Relaying is a duty of every reachable host. The design
is in \code{planning/cpp-relay.md}; the test \code{host\_mesh\_discovery}
runs a host with \code{--no-listen} that joins through another and
receives a paglet.

\subsection{The channel between hosts}\label{sec:communication-channel}

Hosts talk HTTPS, and inside it Noise channels~\cite{perrin2018noise}
(\code{Noise\_XX\_25519\_ChaChaPoly\_SHA256}, built on libsodium). The
static keys of the handshake are the X25519 forms of the Ed25519 identity
keys, so the peer of a channel is an identity key, which the node checks
against the ledger (Section~\ref{sec:mesh}). The prologue contains the
mesh ID, so hosts of different meshes fail at the second handshake
message. TLS keeps proxies and firewalls working and authenticates
nothing by itself: a host checks the certificates of the hosts it
connects to only when it was started with \code{-{}-tls-ca FILE}; the
identity of a peer is its Noise key in every case. A server handles the
requests of a session one at a time, so
frames arrive in the order they were sent. Frames that cannot be delivered
are dropped and counted; the protocols above time out: a lookup asks the
records once more and then fails, a move offers the next host of its
ticket or is aborted, and the paglet stays.

\subsection{Measuring the links}\label{sec:communication-latency}

The Latency Map (\code{cpp/examples/latency/}) measures what a request
between paglets on two hosts costs. It clones a probe to every host; each
probe derives an endpoint to itself that allows only \code{ping},
\code{probe.peers} and \code{probe.done} and sends it to the parent, which
hands every probe copies of the others' endpoints narrowed to
\code{ping}. A probe pings its targets one at a time, waits for each
reply before the next ping and keeps the shortest and the mean round
trip per target (Listing~\ref{lst:comm-ping}); a ping without a reply in
5\,s puts the status into the cell. The parent collects one row per
probe with the \code{FanOut} pattern and then tells the probes to end.

% snippet: cpp/examples/latency/latency.cpp#ping [float,caption={A probe of the Latency Map pings one target at a time and records the round trip in the continuation},label={lst:comm-ping}]
\begin{lstlisting}[style=paglets-cpp,float,caption={A probe of the Latency Map pings one target at a time and records the round trip in the continuation \src{cpp/examples/latency/latency.cpp}},label={lst:comm-ping}]
// Pings target `t`, round `round`; one at a time.
void ping(std::size_t t, std::int64_t round) {
    if (t >= targets_.size()) {
        for (auto& x : targets_) {
            if (x.cell.replies > 0) x.cell.avg_ms /= static_cast<double>(x.cell.replies);
            row_.cells.push_back(x.cell);
        }
        return pt::report(parent_, paglets::encode(row_));
    }
    if (round >= rounds_) return ping(t + 1, 0);
    const double sent_at = now_ms();
    paglets::RequestOptions o;
    o.timeout_ms = 5000;
    auto sent = targets_[t].endpoint.request(
        "ping",
        [this, t, round, sent_at](Message& reply) {
            Cell& c = targets_[t].cell;
            if (reply.ok()) {
                const double rtt = now_ms() - sent_at;
                c.avg_ms += rtt;
                c.min_ms = c.replies == 0 ? rtt : std::min(c.min_ms, rtt);
                ++c.replies;
            } else {
                c.error = std::string(paglets::abi::error_name(reply.status()));
            }
            ping(t, round + 1);
        },
        std::move(o));
    if (!sent) {
        targets_[t].cell.error = pt::error_text("ping", sent.error());
        ping(t + 1, 0);
    }
}
\end{lstlisting}
% /snippet

Table~\ref{tab:comm-latency} is the last run on real machines: the live
mesh of Section~\ref{sec:evaluation-setup}, a Mac Studio (\emph{mac}), a
Raspberry Pi~4 (\emph{pi}) and a uConsole built on a Raspberry Pi
Compute Module~5 (\emph{uconsole}) on a gigabit LAN, 20 pings per pair,
with the probes in worker processes. A round trip passes both runtimes,
the worker process boundary on both sides and the encrypted channel
between the hosts; relays would add to it. The evaluation
(Section~\ref{sec:evaluation-mesh}) puts the numbers next to the TCP
baseline and the transfer rates of the same mesh.

% Hand-written from docs/results/live_latency.json (recorded 2026-10-12):
% the `% results: live_latency` marker renders every column of the JSON
% (mode, outcome, build, date, ...), which is wider than the text block;
% see the refinement plan. Keep these numbers in step with the JSON.
\begin{table}[htbp]
\centering\small
\caption{Round trips of a request between probes on two hosts (mac: Mac Studio; pi: Raspberry Pi~4; uconsole: uConsole on a Raspberry Pi Compute Module~5), gigabit LAN, probes in worker processes, Release build, 20 pings per pair, recorded 2026-10-12.}
\label{tab:comm-latency}
\begin{tabular}{lrr}
\toprule
Pair & Mean of 20 (ms) & Fastest (ms) \\
\midrule
mac $\rightarrow$ pi & 2.84 & 1.08 \\
mac $\rightarrow$ uconsole & 4.69 & 1.92 \\
pi $\rightarrow$ mac & 4.07 & 1.71 \\
pi $\rightarrow$ uconsole & 1.34 & 0.79 \\
uconsole $\rightarrow$ mac & 5.16 & 1.83 \\
uconsole $\rightarrow$ pi & 1.16 & 0.79 \\
\bottomrule
\end{tabular}
\end{table}
